%=====================================================================
\chapter{Quicksort, Randomisation and the Average Case}\label{ch:quicksort}
%=====================================================================
\locator{2}{Part II --- Recursion and Divide-and-Conquer}

\begin{outcomes}
By the end of this chapter you will be able to:
\begin{tight}
  \item \textbf{Trace} Lomuto partitioning and \textbf{state} its invariant.
  \item \textbf{Derive} quicksort's $\bigTh{n^{2}}$ worst case and
        $\bigTh{n\lg n}$ expected case, and \textbf{identify} the input that
        causes each \emph{(CLO-1, CLO-2)}.
  \item \textbf{Explain} why randomising the pivot converts an assumption about
        the data into an assumption about a random number generator.
  \item \textbf{Account} for quicksort's dominance in practice despite a worse
        worst case than merge sort.
  \item \textbf{Predict} the recursion depth of a sort and \textbf{relate} it
        to the stack a program is given.
\end{tight}
\end{outcomes}

\begin{prereq}
\chref{cases} for the distinction between an average over inputs and an
expectation over an algorithm's own coins --- this chapter is that distinction
made concrete. \chref{mergesort} for the pattern and the baseline.
\end{prereq}

\begin{keyterms}
quicksort $\bullet$ partition $\bullet$ Lomuto $\bullet$ pivot $\bullet$
in-place $\bullet$ unbalanced partition $\bullet$ expected running time
$\bullet$ randomised algorithm $\bullet$ adversary $\bullet$ median-of-three
$\bullet$ introsort $\bullet$ tail recursion $\bullet$ recursion depth
\end{keyterms}

\section{Hard Split, Easy Join}

Quicksort is merge sort's mirror image. It does all its work in the divide
step and nothing in the combine step.

\dsfig{\aoaalg{\algname{Quicksort}$(A, \mathit{lo}, \mathit{hi})$}{
\If{$\mathit{hi} - \mathit{lo} < 2$}
  \State \Return
\EndIf
\State $p \gets$ \algname{Partition}$(A, \mathit{lo}, \mathit{hi})$
  \Comment{all the work is here}
\State \algname{Quicksort}$(A, \mathit{lo}, p)$
\State \algname{Quicksort}$(A, p+1, \mathit{hi})$
  \Comment{no combine step at all}
}}{\algname{Quicksort}. There is no line after the two recursive calls: once
the pieces are sorted, the array is sorted, because \algname{Partition} already
put them in the right places.}{quicksort}

\dsfig{\aoaalg{\algname{Partition}$(A, \mathit{lo}, \mathit{hi})$}{
\State $x \gets A[\mathit{hi}-1]$
  \Comment{the pivot: the last element}
\State $i \gets \mathit{lo}$
\For{$j \gets \mathit{lo}$ \textbf{to} $\mathit{hi}-2$}
  \If{$A[j] \le x$}
    \State exchange $A[i] \leftrightarrow A[j]$
    \State $i \gets i + 1$
  \EndIf
\EndFor
\State exchange $A[i] \leftrightarrow A[\mathit{hi}-1]$
\State \Return $i$
}}{\algname{Partition} (Lomuto). Everything $\le$ the pivot is swept to the
left; the pivot is then dropped into the gap. Linear, in place, and the only
place quicksort compares anything.}{partition}

\begin{theorem}[Partition is correct]\label{thm:partition}
After \algname{Partition} returns $i$, every element of
$A[\mathit{lo}\mathinner{\ldotp\ldotp}i]$ is $\le A[i]$, every element of
$A[i+1\mathinner{\ldotp\ldotp}\mathit{hi}]$ is $> A[i]$, and $A[i]$ is the
original pivot.
\end{theorem}

\begin{proof}
Invariant for the \textbf{for} loop: \emph{every element of
$A[\mathit{lo}\mathinner{\ldotp\ldotp}i]$ is $\le x$; every element of
$A[i\mathinner{\ldotp\ldotp}j]$ is $> x$; and $A[\mathit{hi}-1] = x$.}

\smallskip
\textbf{Initialisation.} $i = j = \mathit{lo}$, so both ranges are empty and
both clauses hold vacuously; the pivot has not been moved.

\smallskip
\textbf{Maintenance.} If $A[j] > x$, the loop body does nothing and $j$
advances, extending the second range by an element that is indeed $> x$. If
$A[j] \le x$, the body exchanges $A[i]$ with $A[j]$: position $i$ --- which
held an element $> x$, or nothing if the range was empty --- receives an
element $\le x$, extending the first range, and the displaced element lands at
$j$, which is now inside the second range and is $> x$. Incrementing $i$ and
$j$ restores the invariant.

\smallskip
\textbf{Termination.} $j$ increases by one each iteration towards the fixed
bound $\mathit{hi}-2$, so the loop ends. At that point $A[\mathit{lo}
\mathinner{\ldotp\ldotp}i]$ holds elements $\le x$ and $A[i\mathinner{\ldotp
\ldotp}\mathit{hi}-1]$ elements $> x$. Line 8 exchanges the pivot into position
$i$; everything left of $i$ is still $\le x$ and everything right of it is
$> x$, which is the claim.
\end{proof}

\section{Two Running Times, Both True}

\begin{theorem}[Quicksort's worst case]\label{thm:qsworst}
\algname{Quicksort} performs $\bigTh{n^{2}}$ comparisons in the worst case, and
that worst case occurs on already-sorted input.
\end{theorem}

\begin{proof}
If at every level the pivot is the largest (or smallest) remaining element, the
partition splits $n$ elements into $n-1$ and $0$, giving
\[ T(n) = T(n-1) + T(0) + \bigTh{n} = T(n-1) + \bigTh{n}. \]
Expanding, $T(n) = \sum_{k=1}^{n} \bigTh{k} = \bigTh{n^{2}}$ --- the recursion
tree is a path of length $n$, not a balanced tree of depth $\lg n$.

\smallskip
That this is attained by sorted input is immediate: the pivot is $A[\mathit{hi}
-1]$, the last element, which in a sorted array is the largest. Every element
satisfies $A[j] \le x$, so $i$ advances every iteration and the final exchange
returns $i = \mathit{hi}-1$. The split is $(n-1, 0)$ at every level.
\end{proof}

\begin{alertbox}[The worst case is the commonest input shape in practice]
This is the part that matters. Quicksort's worst case is not an exotic
adversarial construction --- it is \emph{data that is already in order}, which
is what you get from a database query with an \texttt{ORDER BY}, a log file, a
sorted import, or any list a user has already sorted once.

\smallskip
An algorithm whose worst case is a rare accident is tolerable. An algorithm
whose worst case is Tuesday is not.
\end{alertbox}

\begin{theorem}[Randomised quicksort's expected time]\label{thm:qsexpected}
If the pivot is chosen uniformly at random from the subarray, \algname{Quicksort}
performs $\bigO{n \lg n}$ comparisons in expectation, on \emph{every} input.
\end{theorem}

\begin{proof}
Let $z_{1} < z_{2} < \cdots < z_{n}$ be the elements in sorted order and let
$X_{ij}$ indicate that $z_{i}$ and $z_{j}$ are ever compared. Two elements are
compared at most once, since a comparison happens only against a pivot and the
pivot is then excluded from both subproblems. So the total is $\sum_{i<j}
X_{ij}$ and, by linearity of expectation,
$\mathbb{E}[\text{comparisons}] = \sum_{i<j}\Pr[z_{i} \text{ and } z_{j}
\text{ compared}]$.

\smallskip
Consider the set $Z_{ij} = \{z_{i}, \ldots, z_{j}\}$, of size $j-i+1$. While
all of $Z_{ij}$ sits in one subproblem, no element of it has been a pivot. The
first pivot chosen from $Z_{ij}$ decides matters: if it is $z_{i}$ or $z_{j}$,
those two are compared; if it is any of the $j-i-1$ elements strictly between,
then $z_{i}$ and $z_{j}$ are separated into different subproblems and are never
compared. Each element of $Z_{ij}$ is equally likely to be the first pivot
chosen from it, so
\[ \Pr[z_{i} \text{ and } z_{j} \text{ compared}] = \frac{2}{j-i+1}. \]
Therefore
\[ \mathbb{E}[\text{comparisons}]
   = \sum_{i=1}^{n-1}\sum_{j=i+1}^{n} \frac{2}{j-i+1}
   < \sum_{i=1}^{n-1} \sum_{k=1}^{n} \frac{2}{k}
   = 2(n-1)H_{n} = \bigO{n \lg n}, \]
using $H_{n} = \bigTh{\lg n}$ for the harmonic number. The expectation is over
the pivot choices only, with the input fixed, so the bound holds for every
input.
\end{proof}

\begin{conceptbox}[What randomising actually bought]
Compare the two statements carefully.

\smallskip
\textbf{Without randomisation:} ``$\bigTh{n\lg n}$ \emph{on average, assuming
the input is a uniform random permutation}.'' That assumption is about the
world, and the world violates it constantly --- sorted data is common, not
rare.

\smallskip
\textbf{With randomisation:} ``$\bigO{n\lg n}$ \emph{expected, on every
input}.'' No assumption about the data at all. An adversary who has read your
source code, knows your algorithm and chooses the worst possible input still
cannot make it slow, because he does not know your coin flips.

\smallskip
This is \chref{cases}'s distinction, and it is the single most valuable idea in
the chapter: \textbf{randomisation converts an assumption you cannot enforce
into one you can}.
\end{conceptbox}

\section{Measured: the Gap, and Closing It}

\begin{worked}[Comparisons on Already-Sorted Input]
\texttt{code/ch09\_quick.cpp} counts comparisons exactly --- no timing, no
noise.

\rowcolors{2}{white}{lightbg}
\begin{center}
\begin{tabular}{@{}rrrrrr@{}}
\toprule
\rowcolor{primary!15}
$n$ & last-pivot & $/(n^{2}/2)$ & random pivot & $/(n\lg n)$ & speed-up \\
\midrule
 1{,}000 &        499{,}500 & 0.9990 &      12{,}001 & 1.2042 &  41.6 \\
 2{,}000 &      1{,}999{,}000 & 0.9995 &      25{,}000 & 1.1399 &  80.0 \\
 4{,}000 &      7{,}998{,}000 & 0.9998 &      52{,}511 & 1.0971 & 152.3 \\
 8{,}000 &     31{,}996{,}000 & 0.9999 &    121{,}896 & 1.1752 & 262.5 \\
16{,}000 &    127{,}992{,}000 & 0.9999 &    271{,}253 & 1.2139 & 471.9 \\
\bottomrule
\end{tabular}
\end{center}

\smallskip
\textbf{The last-pivot column is $n(n-1)/2$ exactly} --- $0.9990 \to 0.9999$,
converging to 1. Theorem~\ref{thm:qsworst} predicted the split $(n-1,0)$ at
every level, giving $\sum_{k=1}^{n-1}k = n(n-1)/2$ comparisons, and that is
precisely what was counted. This is not an asymptotic estimate; it is the exact
number.

\smallskip
\textbf{The random-pivot column is about $1.15\,n\lg n$}, as
Theorem~\ref{thm:qsexpected}'s $2n\ln n = 1.39\,n\lg n$ bound permits.

\smallskip
\textbf{And the speed-up grows without limit}: 41.6 at $n=1{,}000$ and 471.9 at
$n=16{,}000$. It is multiplying by about 1.8 per doubling, which is what
$n^{2}/(n\lg n) = n/\lg n$ does. Extrapolated to a million elements the factor
is about 30{,}000.
\end{worked}

\begin{worked}[The Worst Case Does Not Merely Run Slowly]
Recursion depth on the same sorted input:

\rowcolors{2}{white}{lightbg}
\begin{center}
\begin{tabular}{@{}rrrr@{}}
\toprule
\rowcolor{primary!15}
$n$ & last-pivot & random pivot & $\lg n$ \\
\midrule
 1{,}000 &    999 & 21 & 10.0 \\
 2{,}000 &  1{,}999 & 26 & 11.0 \\
 4{,}000 &  3{,}999 & 27 & 12.0 \\
 8{,}000 &  7{,}999 & 30 & 13.0 \\
16{,}000 & 15{,}999 & 30 & 14.0 \\
\bottomrule
\end{tabular}
\end{center}

\smallskip
\textbf{The naive version's depth is exactly $n-1$.} Each of those is a live
stack frame. The program measured its own frame size by comparing the address
of a local at depth 0 with one at the deepest depth:

\begin{center}
\begin{tabular}{@{}lr@{}}
depth reached       & 15{,}999 \\
stack span          & 1{,}023{,}912 bytes \\
bytes per frame     & 64.0 \\
default stack       & 1{,}048{,}576 bytes (1 MB) \\
frames that fit     & 16{,}384 \\
\end{tabular}
\end{center}

\smallskip
\textbf{At $n = 16{,}000$ this program used 99.9\% of its stack.} It was 385
frames --- six elements of input --- from dying.

\smallskip
\textbf{And at $n = 32{,}000$ it did die}, with \texttt{0xC00000FD},
\texttt{STATUS\_STACK\_OVERFLOW}, which is why the timing table in the lab
stops at 16{,}000. The prediction ``$n > 16{,}384$ crashes'' and the
observation agree.

\smallskip
\textbf{This is the practical form of the worst case.} A quadratic sort is a
performance bug you can file. A stack overflow is a crash, with no stack trace
worth reading, on the day a customer's data happened to arrive in order.
\end{worked}

\begin{worked}[And the Trap: On Random Input the Naive Version Wins]
\rowcolors{2}{white}{lightbg}
\begin{center}
\begin{tabular}{@{}rrrr@{}}
\toprule
\rowcolor{primary!15}
$n$ & last-pivot (ms) & random pivot (ms) & ratio \\
\midrule
100{,}000 &  7.39 &  8.86 & 0.83 \\
200{,}000 & 16.33 & 17.77 & 0.92 \\
400{,}000 & 33.82 & 37.04 & 0.91 \\
800{,}000 & 73.40 & 77.72 & 0.94 \\
\bottomrule
\end{tabular}
\end{center}

\smallskip
\textbf{On random input the unsafe version is 6--17\% faster}, because it skips
a random number generation and a swap per partition.

\smallskip
\textbf{That is the trap in its purest form.} A developer benchmarking on
shuffled test data will find randomisation costs 10\% and remove it, and the
benchmark will support the decision every single time --- until the day real,
ordered data arrives and the program crashes.

\smallskip
\textbf{Ten per cent is the insurance premium.} The thing insured against is a
factor of 472 and a stack overflow.
\end{worked}

\section{What Is Actually Shipped}

\begin{conceptbox}[Introsort, and the three defences]
\texttt{std::sort} beat every merge sort in \chref{mergesort} by 39\%. It is
\term{introsort}, and it stacks three ideas:

\begin{tightnum}
  \item \textbf{Quicksort} as the main engine, for its small constant and its
        in-place operation --- it touches half the memory merge sort does,
        which is where most of the 39\% comes from.
  \item \textbf{An insertion-sort cutoff} at around 16--32 elements, for
        exactly the reason \chref{role} measured and \chref{mergesort} tuned.
  \item \textbf{A heapsort fallback.} The recursion depth is tracked, and if it
        exceeds about $2\lg n$ the implementation abandons quicksort and
        finishes with heapsort, which is $\bigO{n\lg n}$ in the worst case.
\end{tightnum}

\smallskip
The third is the interesting one: rather than make the bad case unlikely, it
makes it \emph{survivable}. The guarantee becomes $\bigO{n\lg n}$ worst case,
with quicksort's constants almost all of the time.

\smallskip
\textbf{Median-of-three} --- taking the median of the first, middle and last
elements as pivot --- is a cheaper partial defence, and it fixes the
sorted-input case specifically. It does not give a worst-case guarantee: an
adversary who knows you use it can still construct a quadratic input. Use it
with randomisation, not instead of it.
\end{conceptbox}

\begin{conceptbox}[Merge sort or quicksort?]
\rowcolors{2}{white}{lightbg}
\begin{longtable}{@{}L{3.4cm}L{5.0cm}L{5.6cm}@{}}
\toprule
\rowcolor{primary!15}
& \textbf{Merge sort} & \textbf{Quicksort (randomised)} \\
\midrule
\endhead
worst case & $\bigTh{n\lg n}$ guaranteed & $\bigTh{n^{2}}$, probability
  vanishingly small \\
expected & $\bigTh{n\lg n}$ & $\bigTh{n\lg n}$, smaller constant \\
space & $\bigTh{n}$ auxiliary & $\bigTh{\lg n}$ stack, in place \\
stable & yes & no \\
depth & $\bigTh{\lg n}$ always & $\bigTh{\lg n}$ expected \\
choose it when & stability matters; worst case must be bounded; data is on
  disk or a linked list & memory is tight; average speed matters; data is an
  array in memory \\
\bottomrule
\end{longtable}
\end{conceptbox}

\begin{pitfall}
\begin{tight}
  \item \textbf{Benchmarking only on shuffled data.} It is the one
        distribution on which the dangerous version looks good --- measured,
        10\% better.
  \item \textbf{Believing median-of-three is a guarantee.} It defeats sorted
        input, not an adversary.
  \item \textbf{Forgetting recursion depth.} Depth $n-1$ is a crash, not a slow
        program, and it arrives at 16{,}384 elements on a 1 MB stack.
  \item \textbf{Not recursing on the smaller side first.} Recursing into the
        smaller part and looping on the larger bounds the stack at
        $\bigO{\lg n}$ even in the worst case --- a two-line change that
        removes the crash, though not the quadratic time.
  \item \textbf{Using quicksort when stability is required.} It is not stable,
        and no amount of pivot selection makes it so.
  \item \textbf{Seeding the generator with a constant.} An adversary who can
        read your binary can then reproduce your pivots, and the guarantee of
        Theorem~\ref{thm:qsexpected} is gone.
  \item \textbf{Confusing ``expected $\bigTh{n\lg n}$'' with ``average-case
        $\bigTh{n\lg n}$''.} The first holds on every input; the second assumes
        a distribution.
\end{tight}
\end{pitfall}

\begin{lab}[The Price of Insurance]
Reproduces all five tables. Expect thirty minutes.

\begin{lstlisting}[language=C++]
// ch09_quick.cpp -- quicksort, and why randomising is not superstition.
// Build: cl /O2 /EHsc /std:c++17 ch09_quick.cpp     (see Appendix A)

// --- 1. Lomuto partition, the textbook one -------------------------
// Everything <= the pivot is swept to the left. Returns where the
// pivot ended up. Linear, and where every comparison happens.
static size_t partition_last(std::vector<int>& a, size_t lo, size_t hi) {
    int pivot = a[hi - 1];
    size_t i = lo;
    for (size_t j = lo; j + 1 < hi; ++j) {
        ++comps;
        if (a[j] <= pivot) std::swap(a[i++], a[j]);
    }
    std::swap(a[i], a[hi - 1]);
    return i;
}

// --- 2. The same partition, after moving a random element to the end
// One extra swap. That is the entire cost of the change, and it buys
// an expected-time bound that holds for EVERY input -- no assumption
// about the data at all (Chapter 2's distinction between an average
// over inputs and an expectation over the algorithm's own coins).
static size_t partition_random(std::vector<int>& a, size_t lo, size_t hi) {
    std::uniform_int_distribution<size_t> d(lo, hi - 1);
    std::swap(a[d(rng)], a[hi - 1]);
    return partition_last(a, lo, hi);
}

// --- 3. Comparisons counted: the two bounds, seen directly ---------
// On sorted input the last-element pivot splits n into 0 and n-1
// every single time, so the count is 1+2+...+(n-1) = n(n-1)/2
// exactly. Compare that with the random pivot on the same input.

// --- 4. Recursion depth, which is what actually crashes ------------
// The worst case is not only slow. Depth n means n stack frames.

// --- 5. And the same thing in seconds ------------------------------
// Capped at 16,000 deliberately: see part 6 for why 32,000 is not an
// option on a default stack.

// --- 6. The worst case does not just run slowly; it crashes --------
// Measure how big one frame actually is by looking at where a local
// variable sits at depth 0 and at the deepest depth, then work out
// how many fit in the default 1 MB stack. This is a prediction that
// can be checked, and it was: at n = 32,000 this program died with
// 0xC00000FD, STATUS_STACK_OVERFLOW.
double span = (double)(stack_top - stack_deep);
double per  = span / maxdepth;
printf("   bytes per frame      : %.1f\n", per);
printf("   frames that fit      : %.0f\n", 1048576.0 / per);

// --- 7. On RANDOM input the naive pivot is fine --------------------
// Which is the trap: the defect is invisible until the day someone
// feeds your sort data that is already in order, and that day comes.
\end{lstlisting}

\textbf{What to notice.} Part 6 measures the stack rather than assuming it. The
frame came out at exactly 64.0 bytes, which is not a number you could have
looked up --- it depends on this compiler, these optimisation flags and the
number of live locals.

\smallskip
\textbf{Then try to break it.} Make the recursion sort the \emph{smaller} side
recursively and loop on the larger one. Predict first what happens to (a) the
comparison count on sorted input and (b) the recursion depth. The answer is
that (a) does not change at all and (b) drops to $\bigO{\lg n}$ --- so the
crash disappears and the quadratic time does not. Fixing a symptom is not
fixing a bug, and knowing which you have done is the point.
\end{lab}

\begin{checkpoint}
\begin{tightnum}
  \item On sorted input the measured comparison count was exactly $n(n-1)/2$.
        Derive that from the split Theorem~\ref{thm:qsworst} describes.
  \item Explain why ``expected $\bigO{n\lg n}$ on every input'' is a stronger
        claim than ``average-case $\bigO{n\lg n}$'', and why an adversary makes
        the difference \emph{(CLO-2)}.
  \item The measured frame was 64 bytes and the stack 1 MB. At what $n$ does
        the naive sort on sorted input crash, and what would doubling the stack
        buy you?
\end{tightnum}
\end{checkpoint}

\begin{chaptersummary}
\begin{tight}
  \item Quicksort is merge sort inverted: a hard split (partition, $\bigTh{n}$,
        in place) and no join at all.
  \item Worst case $\bigTh{n^{2}}$, attained on \emph{already-sorted} input ---
        which is a common input shape, not an exotic one. Expected
        $\bigO{n\lg n}$ with a random pivot, on every input.
  \item Randomisation converts an assumption about the data, which the world
        violates, into an assumption about a random number generator, which it
        does not.
  \item Measured: on sorted input the naive pivot made exactly $n(n-1)/2$
        comparisons ($0.9999 \times n^{2}/2$) against the random pivot's
        $1.15\,n\lg n$ --- a speed-up of 41.6 at $n{=}1{,}000$ rising to 471.9
        at $n{=}16{,}000$.
  \item Measured: the naive version's recursion depth is exactly $n-1$. At
        $n = 16{,}000$ it used 99.9\% of a 1 MB stack, at 64.0 bytes per frame;
        at $n = 32{,}000$ it crashed with \texttt{STATUS\_STACK\_OVERFLOW}.
  \item Measured: on \emph{random} input the unsafe version is 6--17\% faster.
        That 10\% is the insurance premium, and the thing insured against is a
        factor of 472 and a crash.
  \item \texttt{std::sort} is introsort: quicksort, plus an insertion-sort
        cutoff, plus a heapsort fallback once the depth exceeds $2\lg n$ ---
        making the bad case survivable rather than merely unlikely.
\end{tight}
\end{chaptersummary}

\begin{reviewq}
  \item \textbf{(MCQ)} Quicksort with a last-element pivot achieves its worst
        case on: (a)~random input; (b)~already-sorted input; (c)~input with
        many duplicates only; (d)~no realistic input.
  \item \textbf{(MCQ)} Randomised quicksort's expected $\bigO{n\lg n}$ bound
        assumes: (a)~the input is a random permutation; (b)~the input has no
        duplicates; (c)~nothing about the input; (d)~the input is partly
        sorted.
  \item \textbf{(Short)} State Lomuto partition's loop invariant and explain
        what the final exchange accomplishes.
  \item \textbf{(Short)} Explain why quicksort is preferred in practice despite
        a worse worst case than merge sort \emph{(CLO-1)}.
  \item \textbf{(Short)} Why is median-of-three a defence against sorted input
        but not against an adversary?
  \item \textbf{(Applied)} In the proof of Theorem~\ref{thm:qsexpected}, explain
        why $z_{i}$ and $z_{j}$ are compared exactly when the first pivot
        chosen from $Z_{ij}$ is one of them, and evaluate the resulting
        probability for $i=3$, $j=10$.
  \item \textbf{(Applied)} Your team benchmarks a sort on shuffled data and
        finds randomisation costs 10\%, so proposes removing it. Write the
        reply, citing the measured numbers in this chapter, and name the one
        experiment that would settle the question.
\end{reviewq}
